\section*{Exercises about Subspaces}

\exercise{1}
\begin{xrcs}
  For each of the following subsets of $\myF^3$, determine whether it is a subspace of $\myF^3$:
  \begin{enumerate}
    \item $\left \{ \left(x_1, x_2, x_3 \right ) \in \myF^3 : x_1 + 2x_2 + 3 x_3 = 0 \right \}$
    \item $\left \{ \left(x_1, x_2, x_3 \right ) \in \myF^3 : x_1 + 2x_2 + 3 x_3 = 4 \right \}$
    \item $\left \{ \left(x_1, x_2, x_3 \right ) \in \myF^3 : x_1 x_2 x_3 = 0 \right \}$
    \item $\left \{ \left(x_1, x_2, x_3 \right ) \in \myF^3 : x_1 = 5x_3 \right \}$
  \end{enumerate}

  \begin{xsol}
    \begin{enumerate}
    \item Let $E := \{(x_1, x_2, x_3) \in \myF^3 \mid x_1 + 2x_2 + 3 x_3 = 0\}$. Then $E$ is a subspace of $\myF^3$ because it is a plane through the origin. The additive identity $0 = (0,0,0) \in E$ since $0 + 2(0) + 3(0) = 0$. For $u,v \in E$ with $u_1 + 2u_2 + 3u_3 = 0$ and $v_1 + 2v_2 + 3v_3 = 0$, we have $(u_1 + v_1) + 2(u_2 + v_2) + 3(u_3 + v_3) = (u_1 + 2u_2 + 3u_3) + (v_1 + 2v_2 + 3v_3) = 0 + 0 = 0$. So $u+v \in E$, and $E$ is closed under addition. If $u \in E$ and $\lambda \in \myF$, then $\lambda u_1 + 2(\lambda u_2) + 3(\lambda u_3) = \lambda (u_1 + 2u_2 + 3u_3) = \lambda \cdot 0 = 0$, so $\lambda u \in E$. Thus, $E$ is closed under scalar multiplication. Therefore, $E$ is a subspace of $\myF^3$ by \ref{thm: conditions for a subspace}.

    \item It is not a subspace because it does not contain the additive identity $(0,0,0)$ (since $0 + 2(0) + 3(0) = 0 \neq 4$). It is also not closed under addition and scalar multiplication.

    \item Let $U := \{(x_1, x_2, x_3) \in \myF^3 \mid x_1 x_2 x_3 = 0\}$. Then $U$ is the union of the $xy$-, $xz$-, and $yz$-planes, because at least one of the coordinates must be zero. However, $U$ is not closed under addition:
    \[
    (1,1,0) \in U \quad \text{and} \quad (0,1,1) \in U, \quad \text{but} \quad (1,1,0) + (0,1,1) = (1,2,1) \notin U.
    \]
    Therefore, $U$ is not a subspace of $\myF^3$.

    \item Let $U := \{(x_1, x_2, x_3) \in \myF^3 \mid x_1 = 5x_3\}$. Since $0 = 5(0)$, $0 \in U$. Let $u,v \in U$. Then $u+v = (5u_3+5v_3, u_2+v_2, u_3+v_3) = (5(u_3+v_3), u_2+v_2, u_3+v_3) \in U$, so $U$ is closed under addition. Similarly, for $\lambda \in \myF$ and $u \in U$, $\lambda u = (5\lambda u_3, \lambda u_2, \lambda u_3) \in U$, so $U$ is closed under scalar multiplication. Hence, $U$ is a subspace of $\myF^3$.
  \end{enumerate}
  \end{xsol}
\end{xrcs}

\exercise{2}
\begin{xrcs}
  Verify all assertions about subspaces in Example 1.35.

  \begin{enumerate}
    \item {
      If $b \in \myF$, then $U := \{(x_1, x_2, x_3, x_4) \in \myF^4 \mid x_3 = 5x_4+b \}$ is a subspace of $\myF^4 \iff b=0$.

      \begin{xsol}
        If $b=0$, the condition $x_3 = 5x_4$ defines a subspace containing $0$ and closed under addition and scalar multiplication. If $b \neq 0$, then $0 = (0,0,0,0) \notin U$ because $0 \neq 5(0) + b = b$.
        Moreover, for $u,v \in U$ we know $u_3 = 5u_4 + b$ and $v_3 = 5v_4 + b$, so the third component of the sum is
        \begin{equation}
          u_3 + v_3 = (5u_4 + b) + (5v_4 + b) = 5(u_4+v_4) + 2b,
        \end{equation}
        and hence
        \begin{equation}
          \begin{aligned}
            u+v &= (u_1+v_1, u_2+v_2, u_3+v_3, u_4+v_4) \\
                &= (u_1+v_1, u_2+v_2, 5(u_4 + v_4) + 2b, u_4+v_4).
          \end{aligned}
        \end{equation}
        For $u+v$ to lie in $U$ again, its third component would have to be $5$ times its fourth component plus $b$.
        The third component equals $5(u_4+v_4) + b$ if and only if $2b = b$, which forces $b=0$.
        Thus, $U$ is a subspace of $\myF^4$ if and only if $b=0$.
      \end{xsol}
    }

    \item{
      The set of continuous real-valued functions on the interval $[0,1]$ is a subspace of $\real^{[0,1]}$.

      \begin{xsol}
        The zero function $f(x) = 0$ is continuous. The sum of two continuous functions is continuous, and any scalar multiple of a continuous function is continuous. Thus, it is a subspace.
      \end{xsol}
    }

    \item{
      The set of differentiable real-valued functions on $\real$ is a subspace of $\real^\real$.

      \begin{xsol}
        We check the three subspace conditions: \\
        \emph{Additive identity:} The constant zero function $f(x) := 0 \quad \forall x \in \real$ is differentiable with $f'(x) = 0 \quad \forall x \in \real$.

        \emph{Closedness under addition:} Let $h(x) := f(x) + g(x)$ for differentiable functions $f, g$. Then $h'(x) = f'(x) + g'(x)$, so $h$ is also differentiable.

        \emph{Closedness under scalar multiplication:} Let $h(x) := c f(x)$ for $c \in \real$. Then $h'(x) = c f'(x)$, so $h$ is also differentiable.
      \end{xsol}
    }

    \item{
      The set $S$ of differentiable real-valued functions $f$ on the interval $(0,3)$ such that $f'(2) = b$ is a subspace of $\real^{(0,3)}$ if and only if $b=0$.

      \begin{xsol}
        For $f(x) := 0$, we have $f'(x) = 0$. The condition $f'(2) = 0 = b$ holds if and only if $b=0$.

        For $h(x) := f(x) + g(x)$, we have $h'(x) = f'(x) + g'(x)$. The condition $h'(2) = f'(2) + g'(2) = 2b = b$ holds if and only if $b=0$.

        For $h(x) := cf(x)$, we have $h'(x) = cf'(x)$. The condition $h'(2) = cf'(2) = cb = b$ holds for all constants $c \in \real$ if and only if $b=0$.
      \end{xsol}
    }
  \end{enumerate}
\end{xrcs}

\exercise{3}
\begin{xrcs}
  Show that the set $S$ of differentiable real-valued functions $f$ on the interval $(-4,4)$ such that $f'(-1) = 3f(2)$ is a subspace of $\real^{(-4,4)}$.

  \begin{xsol}
    We check the three subspace conditions: \\
    \emph{Additive identity:} For the zero function $f(x) := 0$, we have $f'(x) = 0$, so $f'(-1) = 0 = 3(0) = 3f(2)$, meaning $0 \in S$.

    \emph{Closedness under addition:} For $f, g \in S$ and $h(x) := f(x) + g(x)$, $h'(-1) = f'(-1) + g'(-1) = 3f(2) + 3g(2) = 3(f(2)+g(2)) = 3h(2)$. Therefore, $h \in S$.

    \emph{Closedness under scalar multiplication:} For $f \in S, c \in \real$, and $h(x) := cf(x)$, $h'(-1) = cf'(-1) = c(3f(2)) = 3(cf(2)) = 3h(2)$. Therefore, $h \in S$.

    Hence, $S$ is a subspace of $\real^{(-4,4)}$.
  \end{xsol}
\end{xrcs}

\exercise{8}
\begin{xrcs}
  Prove or give a counterexample: If $U$ is a nonempty subset of $\real^2$ such that $U$ is closed under addition and under taking additive inverses (meaning $-u \in U$ whenever $u \in U$), then $U$ is a subspace of $\real^2$.

  \begin{xsol}
    The statement is false. Counterexample: let $U := \integer^2 \subset \real^2$. Then $U$ is closed under addition (the sum of integers is an integer) and closed under taking additive inverses, but $U$ is not closed under scalar multiplication by $\real$. For example, $(1,1) \in U$, but $\frac{1}{2}(1,1) = (\frac{1}{2}, \frac{1}{2}) \notin U$. Therefore, $U$ is not a subspace of $\real^2$.
  \end{xsol}
\end{xrcs}

\exercise{12}
\begin{xrcs}
  Prove that the union $U \cup W$ of two subspaces $U, W$ of $V$ is a subspace of $V \iff U \subseteq W$ or $W \subseteq U$.

  \begin{xprf}
    \Leftarrowdirection Let $U$ and $W$ be subspaces of $V$. If $U \subseteq W$, then $U \cup W = W$, which is a subspace. Similarly, if $W \subseteq U$, then $U \cup W = U$, which is a subspace.

    \Rightarrowdirection Suppose $U \cup W$ is a subspace of $V$. To prove that $U \subseteq W$ or $W \subseteq U$, suppose for the sake of contradiction that $U \not\subseteq W$ and $W \not\subseteq U$.
    Then there exist vectors
    \[
      u \in U \setminus W \quad \text{and} \quad w \in W \setminus U.
    \]
    Since $u, w \in U \cup W$ and $U \cup W$ is a subspace, the sum $u+w$ must belong to $U \cup W$.
    By definition of the union, either $u+w \in U$ or $u+w \in W$:
    \begin{itemize}
      \item \prooffont{Case $u+w \in U$:} Then $w = (u+w) - u \in U$ (since $U$ is a subspace), contradicting $w \notin U$.
      \item \prooffont{Case $u+w \in W$:} Then $u = (u+w) - w \in W$ (since $W$ is a subspace), contradicting $u \notin W$.
    \end{itemize}
    In both cases, we reach a contradiction. Therefore, $U \subseteq W$ or $W \subseteq U$.
  \end{xprf}
\end{xrcs}

\exercise{23}
\begin{xrcs}
  Prove or give a counterexample: If $V_1, V_2, U$ are subspaces of $V$ such that $V = V_1 \oplus U$ and $V = V_2 \oplus U$, then $V_1 = V_2$.

  \begin{xsol}
    The statement is false. Counterexample: let $V := \real^2$ and
    \begin{equation}
      U := \{ (x_1, x_2) \in \real^2 \mid x_2 = 0 \} \quad \text{($x$-axis)}.
    \end{equation}
    Define $V_1$ and $V_2$ by:
    \[
      V_1 := \underbrace{\{(x_1, x_2) \in \real^2 \mid x_1 = 0 \}}_{\text{$y$-axis}} \quad \text{and} \quad
      V_2 := \underbrace{\{ (x_1, x_2) \in \real^2 \mid x_2 = x_1 \}}_{\text{line } y=x}.
    \]
    Then $V_1 \neq V_2$, but $V = \real^2 = V_1 \oplus U = V_2 \oplus U$, since every $(x,y) \in \real^2$ can be uniquely written as:
    \begin{equation}
      (x,y) = \underbrace{(x,0)}_{\in U} + \underbrace{(0,y)}_{\in V_1} \quad \text{and} \quad (x,y) = \underbrace{(x-y,0)}_{\in U} + \underbrace{(y,y)}_{\in V_2}.
    \end{equation}
    Therefore, the statement does not hold in general.
  \end{xsol}
\end{xrcs}

\exercise{24}
\begin{xrcs}
  Suppose $V_1, \ddd, V_m$ are finite-dimensional subspaces of $V$. Prove that $V_1 + \dots + V_m$ is a direct sum if and only if
  \begin{equation}
    \dim(V_1 + \dots + V_m) = \dim V_1 + \dots + \dim V_m.
  \end{equation}
\begin{xprf}
  For each $j \in \{1, \dots, m\}$, choose a basis $v_{j,1}, \dots, v_{j,\dim V_j}$ of $V_j$.
  The concatenated list
  \[
    v_{1,1}, \dots, v_{1,\dim V_1}, \dots, v_{m,1}, \dots, v_{m,\dim V_m}
  \]
  spans $V_1 + \dots + V_m$ and has length $\sum_{j=1}^m \dim V_j$.

  Next we claim that the sum is direct if and only if this concatenated list is linearly independent. This is not a definition and needs an argument, so here it is in both directions.

  Suppose first that the sum is direct, and let
  \begin{equation}
    \sum_{j=1}^{m} \underbrace{\sum_{k=1}^{\dim V_j} c_{j,k} v_{j,k}}_{=: \; u_j \, \in \, V_j} = 0.
  \end{equation}
  Each inner sum $u_j$ lies in $V_j$, so this writes $0$ as $u_1 + \cdots + u_m$ with $u_j \in V_j$. By \ref{thm: condition for a direct sum}, directness forces every $u_j = 0$, and then the linear independence of the basis $v_{j,1}, \ddd, v_{j,\dim V_j}$ of $V_j$ forces every $c_{j,k} = 0$. So the concatenated list is linearly independent.

  Conversely, suppose the concatenated list is linearly independent and $u_1 + \cdots + u_m = 0$ with $u_j \in V_j$. Expanding each $u_j$ in the basis of $V_j$ turns this into a vanishing linear combination of the concatenated list, so all its coefficients are $0$, hence every $u_j = 0$. By \ref{thm: condition for a direct sum} the sum is direct.

  Finally, a spanning list is linearly independent if and only if its length equals the dimension of its span, which holds if and only if
  $\dim(V_1 + \dots + V_m) = \sum_{j=1}^m \dim V_j$.
\end{xprf}
\end{xrcs}